\section{问题分析：Agent 的典型失败模式}

在深入解决方案之前，有必要先理解 Agent 在长时间运行场景下的具体失败表现。Anthropic 团队在实验中观察到了两种主要的失败模式。

\subsection{失败模式一：一次性尝试完成所有任务}

即使使用了 compaction（上下文压缩）技术，前沿模型如 Opus 4.5 在面对高层次的 prompt（如"构建一个 claude.ai 的克隆"）时，仍然倾向于尝试一次性完成整个应用。这种"one-shot"行为导致：

\begin{itemize}
    \item Agent 在实现过程中耗尽 context window
    \item 功能只完成了一半，且没有文档记录
    \item 下一个 session 的 Agent 需要花大量时间猜测之前发生了什么
    \item Agent 反复尝试让基础功能重新运行，浪费大量 token
\end{itemize}

\begin{warningbox}{Compaction 不是万能的}
Compaction 是 Claude Agent SDK 的上下文管理能力，可以在 context window 接近满载时压缩历史信息。但 compaction 并不总能将完全清晰的指令传递给下一个 Agent——被压缩的信息可能丢失关键细节，导致后续 Agent 无法准确还原之前的工作状态。不能仅依赖 compaction 来解决跨 session 的连续性问题。
\end{warningbox}

\subsection{失败模式二：过早宣布任务完成}

第二种失败模式通常出现在项目后期。当一些功能已经被构建出来后，后续的 Agent 实例会浏览项目现状，看到已有一定进展，就直接宣布整个任务完成（"declare the job done"），而实际上还有大量功能未实现。

\subsection{问题的本质分解}

这两种失败模式可以被分解为两个子问题：

\begin{table}[H]
\centering
\begin{tabular}{p{4cm}p{9cm}}
\toprule
\textbf{子问题} & \textbf{描述} \\
\midrule
环境初始化 & 需要为所有待实现的功能搭建基础架构，引导 Agent 按步骤、按功能逐个推进 \\
\midrule
增量式推进 & 每个 Agent session 应该做出增量进展，并在结束时将环境留在一个"干净的状态" \\
\bottomrule
\end{tabular}
\caption{长时间运行 Agent 面临的两个核心子问题}
\end{table}

\begin{importantbox}{什么是"干净的状态"？}
"干净的状态"（clean state）指的是代码达到了可以合入主分支的标准：没有重大 bug，代码组织有序且有良好文档，新的开发者（或新的 Agent session）可以立即开始下一个功能的开发，而不需要先清理之前遗留的混乱。这与人类工程师的最佳实践完全一致。
\end{importantbox}

\subsection{本章小结}

Agent 在长时间运行场景下有两大典型失败模式：一是贪多嚼不烂地试图一次完成所有任务，二是在项目有一定进展后过早宣布完工。这两个问题的根源都在于缺乏结构化的任务分解和进度追踪机制。


